\scanpage{35}
\mspage{16}
On va regarder de plus près la partie $S\hat{\mathcal E}$ de $S\widehat{\widehat{\mathcal E}}$, image inverse de
$\underset{\substack{\Vert\\\hat\Gamma}}{\Gamma}=\unsure{\mathfrak S_3|\mathcal S_3}\times\mathbb Z/2\subset\widehat{\widehat\Gamma}=\operatorname{Autext}_{\mathrm{lac}}\hat\pi$.

On va regarder les éléments $\unsure{\rho}\deleted{\unsure{\tau}},\dot\sigma_0,\sigma_1,\sigma_\infty,\tau_0,\tau_1,\tau_\infty$ de $\widehat{\widehat{\mathcal E}}$, \deleted{\unsure{extérieurs}} (NB. ils sont $\unsure{\bar m|=}$ dans $\mathcal E$) et voir comment ils se relèvent en $S\widehat{\widehat{\mathcal E}}$ (et $\unsure{\bar m|=}$, en $S\mathcal E$). En fait, par $\dot\rho\mapsto\rho,\dot\sigma_0\mapsto\tilde\sigma_0=\sigma_0\tau_0$, on a un relèvement \inserted{multiplicatif} de $\Gamma_P^*=\{(\sigma,\operatorname{sg}(\sigma))\mid\sigma\in\unsure{\mathfrak S_3|\mathcal S_3}\}$ dans $\mathcal E$, qu'on voudrait remonter à $S\mathcal E$ --- et de $\unsure{\bar m|=}$ pour $\dot\sigma_0\mapsto\sigma_0,\tau\mapsto\tau_0$ du ss-groupe de stabilité de $0$ dans $\Gamma$ (etc, par conjugaison, \unsure{idem} pour $1,\infty$). On trouve en fait un homomorphisme canonique
\[
(11)\qquad\varphi:\Gamma_P^*\longrightarrow S\mathcal E
\]
\deleted{qui remonte.} section partielle de $S\mathcal E$ sur $\Gamma$, qui remonte $\Gamma^*\to\mathcal E$ déjà considéré, par
\[
\begin{aligned}
(12)\qquad \dot\rho=(\dot\rho,1)&\longmapsto(1,\rho,1,1,1)\overset{\mathrm{dfn}}{=}\tilde\rho\\
\dot\sigma'_i=(\dot\sigma_i,-1)&\longmapsto(-1,\dot\sigma_i,1,1,1)\overset{\mathrm{dfn}}{=}\tilde\sigma'_i\qquad\&\quad i\in\{0,1,\infty\}.
\end{aligned}
\]

On peut dire aussi que le sous-\deleted{groupe}\inserted{ens.} de $S\widehat{\widehat{\mathcal E}}$ formé des $(p,\sigma,g_0,g_1,g_\infty)$ pour lesquels $g_0,g_1,g_\infty$ sont égaux à $1$ est un \uline{ss-groupe}.
\mspage{34}
